% !TeX root = ../初等数学笔记.tex
% 由主文档通过 \input 载入；请编译主文档。
\setcounter{chapter}{18} % 保留原书章号。
\chapter{平面向量、直线与圆的坐标方法}
\label{chap:em-vectors}
向量表达方向与位移，坐标把这些关系化为实数运算。解析几何的关键不在于写出更多方程，而在于选择坐标并证明代数条件与原几何条件等价。

\section{向量、运算与线性组合}
\emph{向量}（vector）具有大小和方向，可由有向线段表示；同向且等长的有向线段表示同一自由向量。零向量大小为零，不指定方向。沿折线依次移动定义向量加法，$\overrightarrow{AB}+\overrightarrow{BC}=\overrightarrow{AC}$；平行四边形法则给出等价表示。

实数乘法 $\lambda\boldsymbol u$ 把长度乘以 $|\lambda|$，正数保持方向、负数反向，零数给出零向量。形如 $\lambda\boldsymbol u+\mu\boldsymbol v$ 的式子称为线性组合。对非零 $\boldsymbol v$，$\boldsymbol u$ 与它共线当且仅当 $\boldsymbol u=\lambda\boldsymbol v$；若允许零向量，代数共线判定仍可成立，但不能据此谈零向量的几何方向。

\section{坐标、数量积与度量}
选择相互垂直的单位向量 $\boldsymbol e_1,\boldsymbol e_2$ 后，每个平面向量唯一表示为 $u_1\boldsymbol e_1+u_2\boldsymbol e_2$，记为 $(u_1,u_2)$。加法、数乘按坐标进行，勾股定理给出 $|\boldsymbol u|=\sqrt{u_1^2+u_2^2}$，两点距离为 $\sqrt{(x_2-x_1)^2+(y_2-y_1)^2}$。

对两个非零向量，\emph{数量积}（dot product）定义为 $\boldsymbol u\cdot\boldsymbol v=|\boldsymbol u||\boldsymbol v|\cos\theta$，其中 $\theta\in[0,\pi]$ 是夹角；含零向量时定义为零。由式 \eqref{eq:em-cosine-rule} 计算 $|\boldsymbol u-\boldsymbol v|^2$，并与坐标展开比较，得到
\[
 \boldsymbol u\cdot\boldsymbol v=u_1v_1+u_2v_2.
\]
因此数量积具有分配律。非零向量垂直当且仅当数量积为零，平行的坐标判定为 $u_1v_2-u_2v_1=0$。沿非零 $\boldsymbol v$ 的有向投影长度为 $(\boldsymbol u\cdot\boldsymbol v)/|\boldsymbol v|$，符号记录投影方向。

\section{分点、直线方程与位置关系}
在线段 $AB$ 内，若 $AP:PB=m:n$，其中 $m,n>0$，则
\[
 P=\frac{nA+mB}{m+n}
\]
按坐标逐项计算，因为 $\overrightarrow{AP}=\frac m{m+n}\overrightarrow{AB}$。中点是 $m=n$ 的情形；三角形重心坐标为三个顶点坐标的平均，代入每条中线的 $2:1$ 分点式即可核验。

过点 $P_0=(x_0,y_0)$、方向向量 $\boldsymbol v=(v_1,v_2)\ne0$ 的直线写为 $P=P_0+t\boldsymbol v$，$t\in\mathbb R$。与其垂直的非零法向量为 $(A,B)$ 时，一般式为 $Ax+By+C=0$，要求 $A^2+B^2>0$。$B\ne0$ 时才可写成斜率式 $y=kx+b$，$B=0$ 对应竖直线。两直线的法向量成比例时须再检查常数项以区分平行与重合，否则联立有唯一交点。

\section{距离、夹角与面积的坐标计算}
设点 $P=(x_0,y_0)$，直线 $\ell:Ax+By+C=0$。令
\[
 H=P-\frac{Ax_0+By_0+C}{A^2+B^2}(A,B).
\]
代入直线式得到 $H\in\ell$，且 $\overrightarrow{PH}$ 平行于法向量，故 $H$ 为垂足。于是
\begin{equation}\label{eq:em-point-line}
 d(P,\ell)=\frac{|Ax_0+By_0+C|}{\sqrt{A^2+B^2}}.
\end{equation}
两直线夹角通常取锐角或直角，方向向量 $\boldsymbol u,\boldsymbol v$ 给出 $\cos\theta=|\boldsymbol u\cdot\boldsymbol v|/(|\boldsymbol u||\boldsymbol v|)$，适用于含竖直线的情形。

若三角形从同一顶点出发的两边向量为 $\boldsymbol u,\boldsymbol v$，由 $S=|\boldsymbol u||\boldsymbol v|\sin\theta/2$ 以及二维平方和恒等式可得
\[
 S=\frac12|u_1v_2-u_2v_1|.
\]
绝对值使面积与顶点排列方向无关；表达式为零则三点共线，不能再视为非退化三角形。

\section{圆的方程与直线圆的位置}
圆心 $(a,b)$、半径 $r>0$ 的圆由距离定义得到 $(x-a)^2+(y-b)^2=r^2$。对 $x^2+y^2+Dx+Ey+F=0$ 配方，半径平方为 $(D^2+E^2)/4-F$；它为正才是非退化实圆，为零时仅有一点，为负时没有实点。

直线与圆相离、相切或相交，分别对应圆心到直线的距离大于、等于或小于半径。这也可通过联立方程的实解数验证。若 $P=(x_0,y_0)$ 在圆上，切线法向量为 $(x_0-a,y_0-b)$，故切线式为
\[
 (x_0-a)(x-a)+(y_0-b)(y-b)=r^2.
\]
点必须在圆上，才能把该式解释为所给点处的切线。

\section{例题、习题与学习检查}
\begin{example}
求直线 $3x+4y=10$ 与圆 $x^2+y^2=4$ 的位置关系和切点。
\end{example}
\begin{solve}
圆心到直线的距离为 $10/5=2$，恰等于半径，所以相切。用垂足公式或沿法向量计算，切点为 $(6/5,8/5)$；其坐标同时满足两个方程。
\end{solve}
\begin{enumerate}
\item 求过 $(1,2)$ 且法向量为 $(2,-1)$ 的直线。答：$2x-y=0$。
\item 三角形顶点为 $(0,0),(4,0),(1,3)$，求面积和重心。答：$6$ 与 $(5/3,1)$。
\item 判断 $x^2+y^2+2x-4y+6=0$ 是否有实点。答：无，配方得 $(x+1)^2+(y-2)^2=-1$。
\end{enumerate}

